% Rise, run, ruin, rejuvenation --- plan \S3.3, "\S5".
% Sources: 09/06.  Cut: the turbulence-shedding passage at 09/06:57-79,
% reduced to two sentences carrying no factual assertion, and the Fisher
% self-rebuke box at 09/06:81-105, cut entirely.
% Renames beyond plan \S7.2, all forced by \S7.1: potential N(t) -> V(t)
% (three meanings of N within the chapter, and the coupled variant of
% \cref{var:sec:vitality} already calls the same object V); inflow \xi -> \phi
% (plan \S7.2); consumption coefficient \mu -> c (\mu is the death rate);
% revival rate \alpha -> \omega (\alpha already carries two sanctioned
% meanings in this chapter).  See PHASE_A_REPORT.md 9.1.
% The label collision on the key fig:single-figure between 09/06:55 and 09/01:43
% is resolved by the var: namespace.

\section{Rise, run, ruin, rejuvenation}
\label{var:sec:rrrr}

\subsection{The cycle the model is for}
\label{var:sec:rrrrmotivation}

Where the two competition models work inside an arena of fixed capacity, this
one replaces the capacity with a flow, and the flow buys a cycle in four
movements. \emph{Potential} accumulates while it goes unused; replicative
agencies proliferate into the surplus once there is one; their proliferation
consumes the potential until scarcity brings a catastrophe in which the
agencies are lost; and the potential rebuilds until a revival restarts the
wave.

The idea is Prigogine's. In a system through which energy or matter flows,
structures can emerge that appear to run against the second
law: order appears locally, as a spontaneous reduction of entropy, sustained
by the throughput and destroyed if the throughput stops
\cite{prigogine2018order}. Such \emph{dissipative structures} are
characterised by chaotic dynamics, in which small fluctuations are amplified
into the details of the eventual form. Whether the fragmentation of a genus
into species is an instance of one has been argued
\cite{elitzur1994let,sharma2023assembly}, and the model below does not settle
it; the claim made here is the weaker one, that the four-part cycle is a
recognisable shape worth writing down as a Markov chain. The pair $(X_t,V(t))$ is a birth--death--catastrophe process with
catastrophe rate $\delta X_t/V(t)$ and an auxiliary ODE for $V$, rather than
a new class of object.

\subsection{The model}
\label{var:sec:rrrrmodel}

Let $X_t$ count replicative agencies and $V(t)$ the accumulated potential.
Potential flows in at a constant rate $\phi$ and is consumed at a rate $c$
per agency:
\begin{equation}
\ddt{V}=\phi-c\,X_t .
\label{var:eq:rrrrV}
\end{equation}
Given $V$, the count $X_t$ moves by four channels, of which the second is the
one that carries the model:
\begin{align}
\pr{\Delta X_t=1\mid X_t}&=\lambda X_t\,\Delta t+o(\Delta t)
&&\text{proliferation,}\notag\\
\pr{X_{t+\Delta t}=0\mid X_t}&=\delta\,\frac{X_t}{V(t)}\,\Delta t+o(\Delta t)
&&\text{catastrophe,}\notag\\
\pr{X_{t+\Delta t}=1\mid X_t=0}&=\omega\,\Delta t+o(\Delta t)
&&\text{revival,}\label{var:eq:rrrr}\\
\pr{X_{t+\Delta t}=X_t}&=1-\lrb{\lambda X_t
+\delta\frac{X_t}{V(t)}+\omega\ind\{X_t=0\}}\Delta t+o(\Delta t)
&&\text{no event.}\notag
\end{align}
The catastrophe channel is the model's whole content. Its rate is not
$\delta X_t$, as it would be in the birth--death--catastrophe process of
\ChCore{}, but $\delta X_t/V(t)$, so that catastrophe is driven by per-capita
scarcity: a large population sitting on a large potential is no more at risk
than a small one on a small potential, while a population that has consumed
its potential is at risk in proportion to how thoroughly it has done so. That
is what makes the cycle close rather than run once. Revival at rate $\omega$
out of the empty state keeps $0$ from being absorbing, and is what makes the
process recurrent.

The potential cannot in fact be exhausted, which is what makes the model well
posed. As $V\downarrow0$ with $X_t\ge1$ the catastrophe hazard
$\delta X_t/V(t)$ behaves like $\delta/\bigl(c(t^{\ast}-t)\bigr)$, where
$t^{\ast}$ is the time \cref{var:eq:rrrrV} would reach zero; that integrates
to infinity, so catastrophe arrives almost surely first and $V$ stays positive
for all time.

Neither nonlinearity here is available to the machinery of the preceding
chapters. The catastrophe rate depends on $V$, and $V$ depends on the whole
history of $X$ through \cref{var:eq:rrrrV}, so that the pair $(X_t,V(t))$ is
Markov while $X_t$ alone is not and no generating function for $X_t$
closes.

\subsection{What the simulations show}
\label{var:sec:rrrrsims}

\Cref{var:fig:vbc-two,var:fig:vbc-many} are exact realisations. They show the
cycle running, and they show that its period is not a constant.

\begin{figure}[!htbp]
\centering
\includegraphics[width=0.96\textwidth]{fig07_rrrr_cycles/fig07a}
\caption{Two realisations of \cref{var:eq:rrrrV,var:eq:rrrr}: potential
$V(t)$ in vermillion, agency count $X_t$ in blue
($\phi=4$, $c=0.8$, $\lambda=0.5$, $\delta=0.25$, $\omega=0.15$).
Every cycle has the same four parts in the same order --- potential builds
while $X_t$ is small, $X_t$ grows into the surplus, the surplus is consumed,
the collapse follows --- and no two are copies of one another.}
\label{var:fig:vbc-two}
\end{figure}

\begin{figure}[!htbp]
\centering
\includegraphics[width=0.72\textwidth]{fig07_rrrr_cycles/fig07b}
\caption{Several realisations from identical parameters and identical initial
conditions. The first cycle is nearly common to all of them; later cycles are
not, each catastrophe having reset the phase by a random amount. The
parameters fix the amplitude, not the timing.}
\label{var:fig:vbc-many}
\end{figure}

What $X_t$ counts is deliberately left open. It may be individuals in a
population, species within a genus, or autonomous units within an
organisation, and the model does not distinguish them: it is a proxy for
order, or for locally reduced entropy, whose units are whatever the
application supplies. That generality is also the model's principal weakness.
\Cref{var:sec:ca} is the same hazard committed to one reading of $X$: there
the units are types within a species, the potential is made by the population
rather than supplied to it, and the catastrophe divides instead of
destroying.

\paragraph{Open questions.}\label{var:rem:open-rrrr}
\begin{enumerate}[leftmargin=1.6em,itemsep=0.2em,topsep=0.2em]
\item \textbf{The cycle period.} The distribution of the interval between
successive catastrophes is a first-passage problem for the pair $(X_t,V(t))$,
and is not known even in its mean. The figures establish that the period has
substantial spread; how much, and how it depends on $\phi$, $c$, $\lambda$
and $\delta$, is the quantity to go after.
\item \textbf{A quasi-stationary law.} Conditional on no catastrophe having
occurred, does $X_t$ have a quasi-stationary law? For the constant-rate
birth--death--catastrophe process the answer is known and geometric
(\cref{dist:sec:qsd,dist:thm:qsd}), and the general theory of extinction
times for catastrophe processes is available
\cite{brockwell1985extinction,cairns2004extinction}. Whether the
state-dependent rate $\delta X_t/V(t)$ preserves that structure is the
natural first question to ask of this model, and the one most likely to have
an answer.
\item \textbf{What it could be fitted to.} The model has never been fitted to
anything, and what data would identify $\phi$, $c$ and $\delta$ separately is
not obvious. A time series of $X_t$ alone almost certainly would not, $V$
being unobserved and entering only through the ratio; the question is what
second observable would suffice.
\end{enumerate}

\FloatBarrier
